\documentclass[UTF8,openany,zihao=-4,fontset=none]{ctexbook}
\newcommand{\ManualTitle}{领导监管端软件使用手册}
\newcommand{\ManualShortTitle}{领导监管端软件使用手册｜十量 V3}
\usepackage[a4paper,top=25mm,bottom=24mm,left=25mm,right=22mm,headheight=15pt]{geometry}
\usepackage{fontspec}
\usepackage{xeCJK}
\setmainfont{DejaVu Serif}
\setsansfont{DejaVu Sans}
\setmonofont{DejaVu Sans Mono}
\setCJKmainfont[AutoFakeBold=2.0]{Droid Sans Fallback}
\setCJKsansfont[AutoFakeBold=2.0]{Droid Sans Fallback}
\setCJKmonofont{Droid Sans Fallback}

\usepackage{booktabs}
\usepackage{amssymb}
\usepackage{longtable}
\usepackage{tabularx}
\usepackage{array}
\usepackage{enumitem}
\usepackage{xcolor}
\usepackage{fancyhdr}
\usepackage{hyperref}
\usepackage[most]{tcolorbox}
\usepackage{titlesec}
\usepackage{lastpage}
\usepackage{microtype}

\definecolor{ManualBlue}{HTML}{17365D}
\definecolor{ManualTeal}{HTML}{167D8D}
\definecolor{ManualLight}{HTML}{EDF4F7}
\definecolor{ManualOrange}{HTML}{B85C00}
\definecolor{ManualRed}{HTML}{A61B1B}
\definecolor{ManualGray}{HTML}{53606D}
\definecolor{ManualLine}{HTML}{C9D5DD}

\hypersetup{
  colorlinks=true,
  linkcolor=ManualBlue,
  urlcolor=ManualTeal,
  pdftitle={\ManualTitle},
  pdfauthor={MineGuard 项目组},
  pdfsubject={软件使用手册},
  pdfkeywords={MineGuard, 煤炭, 监管, 企业填报, 使用手册},
  bookmarksnumbered=true,
  bookmarksopen=true
}

\ctexset{
  chapter={
    format=\Huge\bfseries\color{ManualBlue},
    name={第,章},
    number=\chinese{chapter},
    beforeskip=0pt,
    afterskip=22pt
  },
  section={format=\Large\bfseries\color{ManualBlue}},
  subsection={format=\large\bfseries\color{ManualTeal}}
}
\titleformat{\paragraph}[runin]{\bfseries\color{ManualBlue}}{}{0pt}{}[。]

\setlength{\parindent}{2em}
\setlength{\parskip}{0.35em}
\setlist{itemsep=0.25em,topsep=0.35em,leftmargin=2.2em}
\setlist[enumerate,1]{label=\arabic*.}
\renewcommand{\arraystretch}{1.35}
\newcolumntype{Y}{>{\raggedright\arraybackslash}X}
\newcolumntype{P}[1]{>{\raggedright\arraybackslash}p{#1}}

\pagestyle{fancy}
\fancyhf{}
\fancyhead[L]{\small\color{ManualGray}\ManualShortTitle}
\fancyhead[R]{\small\color{ManualGray}文档版本 V1.0}
\fancyfoot[L]{\small\color{ManualGray}内部使用｜技术辅助系统}
\fancyfoot[R]{\small\color{ManualGray}第 \thepage\ 页 / 共 \pageref*{LastPage} 页}
\renewcommand{\headrulewidth}{0.4pt}
\renewcommand{\footrulewidth}{0.4pt}

\fancypagestyle{plain}{
  \fancyhf{}
  \fancyfoot[L]{\small\color{ManualGray}内部使用｜技术辅助系统}
  \fancyfoot[R]{\small\color{ManualGray}第 \thepage\ 页 / 共 \pageref*{LastPage} 页}
  \renewcommand{\headrulewidth}{0pt}
  \renewcommand{\footrulewidth}{0.4pt}
}

\newtcolorbox{importantbox}{
  enhanced,
  breakable,
  colback=ManualLight,
  colframe=ManualBlue,
  boxrule=0.8pt,
  arc=1.5mm,
  left=3mm,right=3mm,top=2mm,bottom=2mm,
  title={重要说明},
  fonttitle=\bfseries
}
\newtcolorbox{warningbox}{
  enhanced,
  breakable,
  colback=orange!6,
  colframe=ManualOrange,
  boxrule=0.8pt,
  arc=1.5mm,
  left=3mm,right=3mm,top=2mm,bottom=2mm,
  title={注意},
  fonttitle=\bfseries
}
\newtcolorbox{dangerbox}{
  enhanced,
  breakable,
  colback=red!4,
  colframe=ManualRed,
  boxrule=0.8pt,
  arc=1.5mm,
  left=3mm,right=3mm,top=2mm,bottom=2mm,
  title={禁止事项},
  fonttitle=\bfseries
}
\newtcolorbox{tipbox}{
  enhanced,
  breakable,
  colback=green!4,
  colframe=ManualTeal,
  boxrule=0.8pt,
  arc=1.5mm,
  left=3mm,right=3mm,top=2mm,bottom=2mm,
  title={操作提示},
  fonttitle=\bfseries
}

\newcommand{\ui}[1]{\textsf{\bfseries #1}}
\newcommand{\code}[1]{\texttt{\small #1}}
\newcommand{\status}[1]{\textsf{#1}}
\newcommand{\blankline}{\par\noindent\rule{\textwidth}{0.4pt}\par}

\newcommand{\MakeManualCover}[5]{
  \begin{titlepage}
    \thispagestyle{empty}
    \vspace*{18mm}
    {\sffamily\bfseries\fontsize{18pt}{22pt}\selectfont\color{ManualTeal} MineGuard 煤炭智能辅助监管系统\par}
    \vspace{8mm}
    \noindent{\color{ManualBlue}\rule{\textwidth}{2.2pt}}
    \vspace{20mm}
    {\centering
      {\sffamily\bfseries\fontsize{31pt}{39pt}\selectfont\color{ManualBlue} #1\par}
      \vspace{9mm}
      {\sffamily\fontsize{16pt}{22pt}\selectfont\color{ManualGray} #2\par}
    }
    \vfill
    \begin{tcolorbox}[colback=ManualLight,colframe=ManualLine,boxrule=0.5pt,arc=1mm]
      \begin{tabularx}{\linewidth}{P{31mm}Y}
        \textbf{适用软件} & #3\\
        \textbf{文档版本} & V1.0\\
        \textbf{发布日期} & 2026 年 8 月 10 日\\
        \textbf{适用对象} & #4\\
        \textbf{文档属性} & 内部使用，正式软件操作与培训材料\\
      \end{tabularx}
    \end{tcolorbox}
    \vspace{8mm}
    {\small\color{ManualGray}#5\par}
    \vspace{12mm}
    {\centering\small\color{ManualGray}MineGuard 项目组\par}
  \end{titlepage}
}

\newcommand{\DocumentControl}[2]{
  \chapter*{文档控制}
  \addcontentsline{toc}{chapter}{文档控制}
  \begin{longtable}{P{34mm}P{110mm}}
    \toprule
    \textbf{项目} & \textbf{内容}\\
    \midrule
    文档名称 & \ManualTitle\\
    软件版本 & #1\\
    文档版本 & V1.0\\
    发布日期 & 2026 年 8 月 10 日\\
    适用范围 & #2\\
    维护原则 & 软件页面、权限或业务流程发生变化时，应同步修订本手册并重新导出 PDF。\\
    \bottomrule
  \end{longtable}

  \section*{修订记录}
  \begin{longtable}{P{20mm}P{30mm}P{25mm}P{70mm}}
    \toprule
    \textbf{版本} & \textbf{日期} & \textbf{状态} & \textbf{说明}\\
    \midrule
    V1.0 & 2026-08-10 & 十量 V3 正式版 & 按当前十量 V3 实现整理角色权限、标准操作、异常处理、安全边界与部署访问说明。\\
    \bottomrule
  \end{longtable}
}


\begin{document}
\MakeManualCover
  {领导监管端软件使用手册}
  {MineGuard 煤矿十量智能辅助监管系统 V3}
  {MineGuard Platform｜十量监管 V3}
  {监管领导、只读查看人员、业务复核人员和平台系统管理员}
  {本手册面向非技术管理人员，以「三步快速查看、十量分组理解、风险优先核查」为主线。领导业务页只读，系统结果不自动认定违法事实、责任或处罚结论。}

\DocumentControl
  {MineGuard 煤矿十量智能辅助监管系统 V3}
  {当前仓库 \code{platform/} 的政府集中平台、领导只读驾驶舱和十量 V3 交换留痕。}

\begin{importantbox}
领导端只接收和分析已登记煤矿 Agent 发送的数据。业务页面没有替企业新增、修改、删除、
确认报表、改变算法参数或手工关闭风险的接口。领导的核心任务是识别优先核查对象，
明确需要调阅的原始资料，并通过本单位既有工作渠道安排人工复核。
\end{importantbox}

\begin{warningbox}
\status{暂未发现异常}只表示当前已收到、可执行的证据内未达到风险门槛，不等于煤矿
安全、合法或合规；\status{数据不足}不能当作正常；企业提交文字说明也不会自动消除
风险。任何监管结论均须结合原始凭证、现场情况和适用制度由有权人员形成。
\end{warningbox}

\chapter*{领导三分钟速览}
\addcontentsline{toc}{chapter}{领导三分钟速览}

\section*{每天只做三步}

\begin{enumerate}
  \item \textbf{看辖区总览。}先看监管煤矿、本期已报、风险煤矿、数据不足和待企业回复。
  未报送和数据不足先安排补数、核对接口或核源，不得解释为无异常。
  \item \textbf{看重点关注。}风险优先排序后，先处理高严重度、连续出现、涉及多类独立
  证据的线索；企业已回复但风险未解除的事项仍需关注。
  \item \textbf{点进一座矿。}查看十量分组、受影响日期、指标、算法依据和趋势，记录需
  调阅的日报、设备记录、班次、地磅、运单、洗选、销售和发票凭证，再通过本单位既有
  渠道明确责任人和时限。
\end{enumerate}

\section*{页面状态的最短解释}

\begin{longtable}{P{38mm}P{58mm}P{48mm}}
\toprule
\textbf{页面状态} & \textbf{领导应怎样理解} & \textbf{建议动作}\\
\midrule
\status{存在风险} & 可核验关系或多源证据形成技术冲突 & 调原始凭证，确定复核范围和时限。\\
\status{数据不足} & 主字段、必需班次或当前可执行证据不足 & 先补数据和核来源，不能当作正常。\\
\status{暂未发现异常} & 当前证据内未发现达到门槛的冲突 & 可抽查；不能写成安全或合法认定。\\
\status{未报送} & 本期没有可用于当前分析的正式报送 & 核对报送计划、网络和企业端状态。\\
\status{Legacy 五量} & 仅有 V2 历史字段 & 只读参考，不补造十量新增字段。\\
\bottomrule
\end{longtable}

\begin{tipbox}
需要会议展示时点击\ui{大屏展示}；需要检索具体煤矿、逐日趋势或风险明细时回到普通
监管工作台。大屏和普通页面都只读，也都要求先登录。
\end{tipbox}

\tableofcontents
\clearpage

\chapter{平台定位与数据边界}

\section{一矿一 Agent，政府集中监管}

\begin{center}
\begin{tabularx}{0.95\textwidth}{Y c Y}
\textbf{各煤矿独立 Agent} & & \textbf{政府 MineGuard Platform}\\
CSV/工作簿/固定目录/受控连接器 & $\longrightarrow$ & 逐矿验签、幂等接收、追加留痕\\
经办复核、负责人确认、可靠发送 & 十量 V3 HTTPS + 双 HMAC & 唯一算法、辖区汇总、只读驾驶舱\\
原因、证据索引和措施 & $\longrightarrow$ & 回复留痕；必要时按修订报表重算\\
\end{tabularx}
\end{center}

每个企业 Agent 只绑定一个煤矿和经营主体，拥有独立系统身份、数据库和密钥。政府
Platform 集中保存已登记煤矿的报送、算法报告、风险送达、企业回复和重算留痕。两端
不共享数据库、目录或运行时代码，只按十量 V3 合同交换。

\section{平台从哪里取得数据}

政府端不会直接读取企业本地文件夹，也不要求企业把 CSV 复制进政府服务器目录。企业
Agent 经人工复核后，通过 HTTPS 把签名报文发送到政府 V3 接口。固定目录和自动连接器
均属于企业 Agent 的采集能力，只能在企业侧生成待复核草稿。

因此「企业源系统有数据」不等于「政府平台已收到」。领导应以页面本期已报、最新报表
期和数据新鲜度为准，并在缺报时同时核对企业报送流程和网络接入。

\section{五量 V2 只读边界}

当前正式主线是十量 V3。五量 V2 历史报文、签名、哈希和旧算法结果只读保留，用于
迁移核查和审计。页面会把缺少新增业务量的记录标成 Legacy 或字段不足；平台不会用
历史、同类矿或模型补造开采、销售、运输、洗煤、开票等字段，也不会把 V2 历史改造成
V3 正式报送。

\chapter{十量怎样看}

\section{三组十量}

领导端把十量分成 4/5/1 三组，便于在几分钟内确定核查方向：

\begin{longtable}{P{34mm}P{57mm}P{53mm}}
\toprule
\textbf{分组} & \textbf{业务量} & \textbf{优先核对资料}\\
\midrule
安全生产支撑（4 项） & 风量、电量、火工品量、入井人员量 & 通风测点、计量表、火工品领退用、人员定位和班次记录。\\
生产煤流（5 项） & 产量、开采量、销售量、运输量、洗煤量 & 生产日报、工作面计量、地磅/运单、销售交付、入洗原煤记录。\\
经营票据（1 项） & 开票量 & 正常/蓝字发票实物吨数及红票、退票、作废等辅助明细。\\
\bottomrule
\end{longtable}

\section{11 个原子字段}

火工品量必须保留两个不同单位的原子字段，因此十个业务量在底层对应 11 个字段：

\begin{longtable}{P{33mm}P{50mm}P{20mm}P{43mm}}
\toprule
\textbf{业务量} & \textbf{V3 原子字段} & \textbf{单位} & \textbf{领导注意点}\\
\midrule
风量 & \code{ventilation\_m3\_min} & \code{m3/min} & 时间加权平均。\\
电量 & \code{electricity\_kwh} & \code{kWh} & 窗口耗电量。\\
雷管 & \code{detonators\_count} & \code{count} & 与炸药不能相加。\\
炸药 & \code{explosives\_kg} & \code{kg} & 与雷管分别展示。\\
入井人员量 & \code{mine\_entry\_persons} & \code{person} & 实际入井人次。\\
产量 & \code{production\_t} & \code{t} & 企业生产报表产量。\\
开采量 & \code{extraction\_t} & \code{t} & 工作面/采掘计量。\\
销售量 & \code{sales\_t} & \code{t} & 销售交付实物量。\\
运输量 & \code{transport\_t} & \code{t} & 出矿/外运净吨数。\\
洗煤量 & \code{wash\_feed\_t} & \code{t} & 入洗原煤量。\\
开票量 & \code{invoiced\_quantity\_t} & \code{t} & 正常/蓝字发票实物吨数。\\
\bottomrule
\end{longtable}

日报必须覆盖 11 个原子字段。前七个原子字段还要求零点、八点、四点班数据；销售、
运输、洗煤和开票以日报为强制口径，班次可选。空白是缺失，不是零；火工品量任一子项
缺失都不能显示为完整火工品量。

\section{销售、运输和开票不能混为一个量}

销售是销售交付实物量，运输是出矿/外运净吨数，开票是正常/蓝字发票对应实物吨数。
三者可能因在途、签收和跨期开票产生合法时差，不能按同一天强制相等，也不能在煤流中
作为三次出库重复扣减。

开票主量必须非负。红票、退票、作废、折让和退货是企业来源系统中的辅助事件，净额
另算。当前主报文不逐笔承载这些事件；涉及开票的核查应同时要求正常发票与辅助明细，
不得只看一个净数。

\chapter{登录、权限与访问范围}

\section{登录}

使用单位批准的 Edge 或 Chrome 打开内网 HTTPS 地址，输入个人账号和密码。正式账号
只显示被授权的煤矿范围；页面隐藏不是唯一保护，服务端会对每个查询再次校验矿井范围。
登录失效时重新登录，不要把账号、密码或会话令牌写入收藏夹 URL、共享文档或大屏脚本。

本机演示可使用 \code{admin / 123123123}，但只允许回环地址受控展示。默认密码、演示
账号和共享 admin 不得用于正式内网。正式账号、矿井范围、停用和重置由平台管理员在
服务器运维命令中维护，领导业务页没有账号管理入口。

\section{当前 V3 的只读权限}

\begin{longtable}{P{32mm}P{55mm}P{57mm}}
\toprule
\textbf{账号类型} & \textbf{可见范围} & \textbf{当前业务能力}\\
\midrule
\code{viewer} & 明确授权的一矿或多矿 & 只读总览、单矿、风险和交换留痕。\\
\code{reviewer} & 明确授权的一矿或多矿 & 当前 V3 领导端仍为只读。\\
\code{supervisor} & 明确授权的一矿或多矿 & 当前 V3 领导端仍为只读。\\
\code{admin} & 全部登记煤矿 & 系统管理员查看和运维；不用于日常共享。\\
\bottomrule
\end{longtable}

页面不存在替企业改数、删除报送、修改算法、审批案件或手工把风险改成已解除的按钮。
如需要交办、催办或形成监管文书，应使用本单位现有经批准业务系统或制度流程；本平台
只提供可追溯的技术线索和企业交换状态。

\chapter{四个只读视图}

\section{辖区总览}

\subsection{顶部六项指标}

\begin{longtable}{P{36mm}P{64mm}P{44mm}}
\toprule
\textbf{指标} & \textbf{含义} & \textbf{领导动作}\\
\midrule
监管煤矿 & 当前账号授权范围内登记煤矿数 & 核对辖区范围是否正确。\\
本期已报 & 当前统计期已有正式报送的煤矿数及覆盖率 & 先查未报送和长期不新鲜。\\
暂未发现异常 & 当前证据内没有达到风险门槛的煤矿数 & 可抽查，不能作合规认定。\\
风险煤矿 & 当前存在未解除风险线索的煤矿数 & 进入风险优先核查。\\
数据不足 & 不能形成确定判断的煤矿数 & 先补字段、班次或核来源。\\
待企业回复 & 已送达且尚待企业说明/修订的事项 & 核对期限和企业交换状态。\\
\bottomrule
\end{longtable}

页面还显示辖区态势、辖区待关注事项、辖区最新动态和煤矿监管矩阵。矩阵可按名称或
编号搜索，显示最新报表期、数据完整率、算法状态、风险、企业回复、数据新鲜度和趋势。
风险项优先排列；点击一座矿进入单矿研判。

\subsection{风险状态变化}

\begin{itemize}
  \item \status{待企业回复}：政府风险报告已形成，等待企业说明或修订数据；
  \item \status{企业已回复、风险未解除}：企业原因、证据索引和措施已留痕，但算法结论未被覆盖；
  \item \status{复核后风险仍存在}：修订或复核后仍触发原风险；
  \item \status{修订复核已解除}：企业提交新的正式报表，平台按同一算法重算后不再出现原风险。
\end{itemize}

\section{单矿研判}

单矿页依次显示矿井身份、唯一十量监管引擎、风险回执、十量一眼看懂、风险优先趋势、
风险与算法依据以及报送/回复时间线。

\subsection{十量一眼看懂}

每项可能显示\status{已提供}、\status{需关注}、\status{数据不足}或\status{未提供}。
「十量已到 10/10」只表示十个业务量具备必要原子数据，不表示所有来源真实或不存在
风险。V2 历史记录的新增业务量显示未提供，而不是用旧值或零补齐。

领导应先看红色需关注项，再看黄色数据不足项。若页面显示字段已接入但整体仍为数据
不足，说明时间覆盖、必需班次、来源或当前可执行模块仍不满足判断条件。

\subsection{风险优先趋势}

趋势图默认最多显示三个业务量并优先选择当前风险项，可点击指标切换。火工品量中的
雷管和炸药保持独立轨道；不同单位不得相加或按纵轴高度互相比较。每条轨道只表示该项
在当前窗口内的自身高低，恒定值位于中线，缺失处断开，悬停可查看原始数值。

趋势区使用完整可见历史帮助观察连续性，当前状态和风险项则以最新报送期为准。领导
优先关注连续偏离、均值突变、多个相关量同向变化和同时伴随数据完整率下降的情况；
单个点只是一条复核线索。

\subsection{风险与算法依据}

每张 finding 卡片展示事项类型、标题、摘要、矿井、类别、时间、状态和最多若干条证据。
核查时记录受影响日期和原子字段，不要只抄一个风险分。时间线用于确认企业报送、政府
分析、报告送达、企业回复和修订重算的顺序。

\section{风险台账}

风险台账是不可变的只读列表，可按事项类型筛选\status{风险线索}或\status{数据待补}，
也可按状态筛选待企业回复、企业已回复未解除、复核后风险仍存在和修订复核已解除。

企业文字回复不能删除旧 finding，也不能把它改成已解除。旧状态和新状态按追加记录
共同保留，便于回答「何时形成、何时送达、企业何时回复、是否有正式修订重算」。

\section{交换留痕}

交换留痕按「企业报送、政府分析、风险送达、企业回复、必要时重算」展示追加式事件。
可选择近 24 小时、近 7 天、近 30 天、本月或自定义时间，并按煤矿、业务环节和业务/
技术记录范围筛选；点击\ui{应用筛选}后可逐页加载更早记录。页面提供
\ui{导出当前筛选}，导出内容仍受当前账号矿井范围限制。

完整性状态不是业务风险状态。若留痕链校验异常，应停止把该页面作为完整审计依据，
保留时间、筛选条件和日志，交平台管理员处理；不得通过删除记录消除提示。

\chapter{大屏展示}

普通监管页右上角点击\ui{大屏展示}可尝试进入全屏。

也可在已登录会话中打开 \path{/wallboard}。浏览器通常要求用户手势才能进入全屏；
必要时使用 F11 或单位批准的 kiosk 模式。

大屏不显示搜索、导出和业务操作，业务数据约每 10 秒刷新，重点煤矿自动轮播。刷新
失败时会保留最后一次成功数据并提示异常，因此领导必须同时查看「数据截至」和连接
状态，不能把旧画面当作实时状态。

大屏中的十量核心图只表达 4/5/1 三组核验范围和当前整体结论，不代表各指标数值大小，
也不会虚构单项状态。会议中若需解释某一矿、某一日或某一指标，应返回普通单矿页。

\chapter{算法结果怎样理解}

\section{唯一十量引擎}

平台对 V3 报送使用唯一十量监管引擎，主要完成：11 原子字段和日报与班次确定性检查、
运行工况识别、生产煤流和经营量的同期间软关系、本矿正式历史和冻结匿名同类矿证据、
时序漂移与变化点、HiGHS 加权 L1 联合协调以及最小冲突集 MCS。算法输出三态结论和
可追溯 finding，不自动预测所谓「真实产量」。

L1 回答的是「在观测容差和软参考区间下，所有量联合相容需要多大最小加权调整」；
MCS 回答的是「最少放宽哪些观测或参考组后才可能相容」。历史关系是辅助证据，不是
物理定律；同类矿只显示匿名区间和样本数，不向企业泄露其他矿明细。

\section{证据不足优先于硬判}

十量主字段、必需班次或当前可执行模块证据不足时，系统返回
\status{数据不足}，不会把「无法验证」显示为正常或冲突。历史冷启动、来源缺失和
折线断开都应保持其原意，不能用邻日、同类矿或模型补数。

\section{当前尚未覆盖的高级资料}

V3 主报文没有期初/期末库存、洗后产品/矸石/损耗、逐批运输/开票日期和来源依赖域。
因此当前页面不代表已经完成原煤收发存、洗选投入产出、逐批凭证账龄或窗口流网络核验；
这些模块在缺少已发布辅助合同的情况下会明确跳过，而不是猜测。交办核查时仍应调阅
上述资料。

\chapter{领导核查记录建议}

平台当前只读，不在页面内保存交办或审批动作。通过本单位既有渠道交办时，建议至少
写清：

\begin{itemize}
  \item 煤矿名称、矿井编号、报表月份和受影响日期；
  \item 页面结论是风险、数据不足还是未报送；
  \item 涉及的业务量和原子字段、单位、日报或班次口径；
  \item 需要调阅的设备读数、生产日报、班次、火工品领退用、人员定位、地磅、运单、
  洗选、销售或发票凭证；
  \item 如果涉及开票，同时提供正常/蓝字发票和红票、退票、作废、折让、退货辅助明细；
  \item 责任人、回复期限和需要提交的更正 V3 报表或证据索引；
  \item finding/报告/消息编号和当时页面的数据截至时间，避免脱离版本讨论。
\end{itemize}

\begin{warningbox}
不要把风险矿数、finding 数、红点数直接写成违法数量、事故数量或处罚等级。企业解释
只能作为一条留痕；若原数据有误，应要求企业依照四眼流程提交新的正式 V3 报表，由
平台重算后再查看状态。
\end{warningbox}

\chapter{常见问题与处理}

\begin{longtable}{P{47mm}P{96mm}}
\toprule
\textbf{现象} & \textbf{处理方法}\\
\midrule
登录后看不到某座矿 & 当前账号未授权该矿，或该矿尚未登记。联系管理员核对矿井范围；不要改 URL 绕过。\\
页面加载但没有数据 & 先看监管煤矿数、本期已报和数据截至。确认企业是否已通过 HTTP 报送、政府是否登记该 sender、页面筛选和账号范围是否正确。\\
本期已报仍显示数据不足 & 查看单矿十量和 finding，核对 11 个日报字段、必需班次、时间覆盖及来源；已收到不等于完整。\\
十量只有部分项目 & 若为 Legacy V2，这是预期只读表现；若为 V3，要求企业从原始来源补正并重新报送，禁止平台补数。\\
火工品量只有一条轨道 & 雷管或炸药另一原子字段缺失。两者单位不同，不能用已有项代表全部火工品量。\\
企业已回复但仍是风险 & 正常语义。说明不会改变算法；只有新的正式修订报表经同一算法重算后，风险才可能解除。\\
趋势图不同指标高低相近 & 各轨道按自身窗口缩放，单位不同不能横向比较绝对高度；悬停查看原始值。\\
大屏数字不更新 & 查看更新失败提示和数据截至；返回普通页刷新。保留最后成功画面不表示实时。\\
交换留痕完整性异常 & 停止以其作为完整审计依据，记录时间和筛选条件，联系管理员检查数据库、磁盘和审计链。\\
浏览器显示页面程序未启动 & 使用新版 Edge/Chrome，退出 IE 模式后按 Ctrl+F5；仍失败时把截图和控制中心日志交管理员。\\
\bottomrule
\end{longtable}

\chapter{每日与会议检查单}

\section{每日三分钟检查}

\begin{enumerate}
  \item 数据截至时间是否为当前预期；本期未报和数据不足分别有多少；
  \item 新增风险、连续风险、企业已回复未解除和超期待回复是否增加；
  \item 风险优先煤矿的十量分组、最多三项趋势和受影响日期是什么；
  \item 需要调阅哪些原始凭证，责任人和期限是否已在本单位业务流程中明确；
  \item 交换留痕是否存在送达、回复或重算异常。
\end{enumerate}

\section{会议展示前检查}

\begin{itemize}
  \item 使用个人只读账号登录，确认授权矿井范围和数据截至时间；
  \item 确认浏览器不是 IE 模式，大屏自动更新正常；
  \item 对重点矿提前在普通单矿页核对日期、字段、单位和 finding 编号；
  \item 对数据不足、未报送、Legacy 和当前未覆盖高级资料作明确口头说明；
  \item 不在投屏、截图或导出中泄露账号、会话、HMAC、企业无关明细或个人敏感信息；
  \item 会后退出账号，公共终端不得保存密码或保持会话。
\end{itemize}

\begin{importantbox}
领导端的正确用法是「先看覆盖与数据不足，再看风险优先级，最后点进单矿确定核查材料」。
它帮助缩小人工核查范围，但不能替代企业原始凭证、现场核验、监管程序和有权人员判断。
\end{importantbox}

\end{document}
